\section{Preparación de datos}
\label{sec:preparacion}

Esta fase transforma las fuentes crudas descritas en la
\hyperref[sec:datos]{sección~\ref*{sec:datos}} en una tabla lista para
modelar, con el objetivo de \textbf{predecir} el nivel del paralelo y la
brecha oficial--paralelo a horizontes cortos. El énfasis del proyecto está en
la previsión, no en el monitoreo en tiempo real: cada decisión de
preparación se toma pensando en evitar fuga de información y en reproducir el
orden temporal real de los datos.

\subsection{Del JSON crudo a los predictores del modelo}

El camino desde la API hasta los predictores tiene cuatro etapas:

\begin{enumerate}
  \item \textbf{Recolección.} Un colector consulta el endpoint
        \texttt{GET /v1/chart/6m} de la API y guarda cada respuesta como
        snapshot JSON en \texttt{data/raw/}: una lista de puntos muestreados
        cada hora.
  \item \textbf{Normalización.} Se elimina el envoltorio de la API y cada
        punto se valida contra un esquema tipado; el resultado se escribe en
        \texttt{data/normalized/chart\_6m.csv}. Toda la API se reduce a
        5 columnas: \texttt{t} (epoch, una fila por hora), \texttt{bb}/\texttt{bs}
        (blue compra/venta, Bs/USD) y \texttt{ob}/\texttt{os} (oficial
        compra/venta).
  \item \textbf{Ingeniería de características.} Una etapa de indexación deriva
        los puntos medios, el spread, la brecha porcentual, el retorno por hora
        y la volatilidad rodante de 24~h; las columnas \texttt{fwd\_*} son
        retornos futuros (etiquetas de previsión, no predictores). Resultado:
        \texttt{data/processed/market\_hourly.csv} (4\,315 filas
        $\times$ 17 columnas).
  \item \textbf{Predictores del modelo.} Los modelos lineales consumen
        rezagos de esos campos más los canales de la Gaceta:
        \begin{itemize}
          \item \textbf{Variante por hora (37 predictores):} rezagos de brecha
                (1, 3 y 24~h), oficial, volatilidad y retorno (todos con
                \texttt{shift(1)}, sin fuga); canales de la Gaceta contados
                por día (\texttt{n\_docs\_fx}, \texttt{n\_docs\_fx\_alta},
                sobre 509 normas etiquetadas por palabras clave) y
                transmitidos a cada hora del día siguiente;
                dummies de hora (23) y de día de semana (6).
        \end{itemize}
\end{enumerate}

\subsection{Remuestreo a frecuencia fija}

Las fuentes no son uniformes: \texttt{history\_full.csv} registra snapshots
por eventos y los exchanges cambian de periodicidad. Todo se alinea a una
\textbf{frecuencia horaria} (4\,321 filas): se conserva el último precio
observado y se interpolan huecos de hasta 3 períodos. La frecuencia diaria se usa
solo para los canales de la Gaceta, que tienen resolución de publicación.


\subsection{Variable objetivo y conjunto de eventos}

El problema se formula como \textbf{previsión supervisada} a dos horizontes:
el nivel del paralelo (o la brecha) \textbf{a una hora y a un día}, con
intervalos que cuantifiquen la incertidumbre. Para cada horizonte $h$ los
predictores usan solo información disponible en $t-h$, y la línea base naïve
repite el último valor observado en ese origen. Para el Objetivo~2 se cruza la serie horaria con las fechas de
edición de la Gaceta: por cada publicación se toma el precio ancla (12:00~UTC)
y se miden retornos a 5~h, 1, 3 y 5 días contra la variación media de todos los
días de la serie (línea base incondicional).

\subsection{Partición temporal y línea base}

La partición respeta el orden temporal, nunca una división aleatoria: el
ventana de prueba es el \textbf{último corte temporal} (las últimas 168~h), el mismo
para todos los modelos. Se define una
\textbf{línea base naïve} (último valor observado, equivalente al paseo
aleatorio) que todo modelo debe superar fuera de muestra: en tipos de cambio
es un referente notoriamente difícil de batir.
